\documentclass[12pt,a4paper]{article}
\usepackage[utf8]{inputenc}
\usepackage[T1]{fontenc}
\usepackage[brazil]{babel}
\usepackage{indentfirst}
\usepackage[margin=2.5cm]{geometry}

\title{Resenha Cr\'itica: Artigo 10 ---\\
Parallel Changes in Large-Scale Software Development}
\author{Gustavo Pessoa Firmino Duarte}
\date{23 de Abril de 2026}

\begin{document}
\maketitle

\section{Introdu\c{c}\~ao}

O artigo sobre mudan\c{c}as paralelas em desenvolvimento de software de larga escala
aborda um dos problemas mais pr\'aticos e comuns da engenharia de software: como
introduzir mudan\c{c}as significativas em sistemas em produ\c{c}\~ao sem interromper
o servi\c{c}o e sem exigir que consumidores de uma API ou biblioteca atualizem sua
depend\^encia de forma s\'incrona. A t\'ecnica conhecida como \textit{Parallel Change}
--- tamb\'em chamada de \textit{Expand and Contract} --- propõe um processo em
tr\^es fases que permite introduzir a nova vers\~ao de uma interface enquanto a
antiga ainda est\'a ativa, migrar os consumidores gradualmente e, por fim, remover
a vers\~ao legada. Essa abordagem reduz drasticamente o risco de grandes migra\c{c}\~oes
e viabiliza a entrega cont\'inua mesmo em sistemas com m\'ultiplos times e
depend\^encias cruzadas.

\section{A T\'ecnica Expand-and-Contract e seus Desafios}

O ciclo de \textit{Parallel Change} ou \textit{Expand-and-Contract} divide-se em
tr\^es fases distintas. Na fase de \textbf{expans\~ao} (\textit{expand}), a nova
interface, comportamento ou estrutura de dados \'e adicionada ao sistema sem remover
a antiga. O sistema passa a suportar ambas as vers\~oes simultaneamente. Na fase
de \textbf{migra\c{c}\~ao} (\textit{migrate}), os consumidores s\~ao gradualmente
atualizados para usar a nova vers\~ao --- essa fase pode durar dias, sprints ou
meses, dependendo da escala e da quantidade de consumidores. Na fase de
\textbf{contra\c{c}\~ao} (\textit{contract}), a vers\~ao antiga \'e removida, o
sistema \'e limpo e a d\'ivida tempor\'aria de compatibilidade retroativa \'e quitada.

Essa estrat\'egia \'e especialmente poderosa para mudan\c{c}as em esquemas de banco
de dados, APIs REST, eventos de dom\'inio e interfaces de m\'odulos internos. Para
esquemas de banco de dados, por exemplo, \'e poss\'ivel adicionar uma nova coluna
\textit{nullable}, migrar os dados gradualmente, garantir que o novo c\'odigo a
popule, e s\'o ent\~ao tornar a coluna NOT NULL e remover a antiga --- sem nenhum
\textit{downtime} e sem necessidade de janelas de manuten\c{c}\~ao.

Os desafios da t\'ecnica incluem a complexidade de manter dois caminhos de c\'odigo
ativos simultaneamente, a disciplina necess\'aria para efetivamente completar a
fase de contra\c{c}\~ao (muitas equipes completam a expans\~ao mas nunca removem
o c\'odigo legado) e a coordena\c{c}\~ao entre times quando os consumidores s\~ao
equipes diferentes com cadências de release pr\'oprias.

\section{Conex\~ao com a Pr\'atica Profissional}

Durante o desenvolvimento do meu MVP de e-commerce, enfrentei em menor escala o
problema que este artigo descreve: precisei renomear uma coluna no banco de dados
Prisma ap\'os j\'a ter dados em produ\c{c}\~ao (mesmo que fosse um ambiente de
desenvolvimento). A solu\c{c}\~ao improvisada foi criar uma migra\c{c}\~ao que
adicionava a nova coluna, copiava os dados e removia a antiga em uma \'unica
opera\c{c}\~ao --- que teria causado downtime em produ\c{c}\~ao real. A t\'ecnica
de Parallel Change teria sido a abordagem correta: adicionar a nova coluna,
atualizar o c\'odigo para escrever em ambas as colunas, migrar os dados existentes
e, s\'o ent\~ao, remover a coluna antiga.

No contexto do est\'agio na ArcelorMittal, as atualiza\c{c}\~oes de sistemas
internacionais frequentemente exigiam coordena\c{c}\~ao entre equipes em pa\'ises
diferentes, cada uma com janelas de deploy distintas. A id\'eia de uma mudan\c{c}a
que possa ser lan\c{c}ada de forma incremental, sem requerer sincroniza\c{c}\~ao
perfeita entre todos os consumidores, \'e extremamente valiosa nesse contexto ---
e \'e exatamente o problema que a t\'ecnica de mudan\c{c}as paralelas resolve.

\section{Conclus\~ao}

A t\'ecnica de \textit{Parallel Change} representa uma aplica\c{c}\~ao concreta
do princ\'ipio de entrega cont\'inua a mudan\c{c}as estruturais em sistemas em
produ\c{c}\~ao. Ao decompor uma mudan\c{c}a potencialmente disruptiva em uma
sequ\^encia de passos seguros e revers\'iveis, ela reduz o risco e aumenta a
confian\c{c}a da equipe em realizar altera\c{c}\~oes significativas. O artigo
tamb\'em nos lembra de uma disciplina essencial: toda fase de expans\~ao deve ter
uma data planejada para a fase de contra\c{c}\~ao, caso contr\'ario o sistema
acumula camadas de compatibilidade retroativa que se tornam d\'ivida t\'ecnica de
dif\'icil pagamento.

\end{document}
